% SPDX-License-Identifier: CC-BY-4.0
% Copyright (c) 2026 Martin Schröder <info@swedishembedded.com>

\section{Limitations}\label{sec:limitations}

The structure claims show that definitions, constants, tests and messages exist at the pinned
commits; they do not show that the code behaves accordingly. Behaviour is established only by
the experiments, which are pilots: one run each per setting, one generator family, small
corpora, 30 facts per generation, one policy size per run. The training settings of E2 were
chosen on two earlier pilots with other seeds, not on the registered run's own measurement, but
they were chosen by hand and no search over them was made.

The registered E2/E3 criterion (a release with a significant gain, then retention and rollback) was not
met: every candidate of the 0.6B policy passed the improvement and retention checks and was refused by the
anchor check, whose bound of 0.02 over 40 questions (0.025 per question) allows no lost answer. One
post hoc run with a 1.7B policy released its first candidate; its second candidate was refused, so the
alias move after a second release, retention against an earlier release as the gate measures it and the rollback
to a first release are described from the code (\cref{sec:gate}) and are not demonstrated end to end here;
the one rollback attempted was refused because the release had no predecessor. In one run the serve check
could not start its server and is not measured. The gate's retention check held the champion to held-out tasks that,
for a facts campaign, are facts it never taught, so it passed while the candidate forgot taught facts
(\cref{sec:e2}); we measured that forgetting on independent probes. The 1.7B result is a single post hoc run.
We did not test a larger anchor suite or other settings that might keep the anchor at 0.6B beyond the two
reported. The facts are numbers, which may be why they leak into numeric answers elsewhere; a
domain of non-numeric facts is not measured.

The held-out families of a run are facts the policy is never taught, a few of the 30 facts of a
generation, and the teacher's answer was not verified for a further share of them
(\cref{sec:defects}); the facts measured in the probes include both, so the counts of facts answered
are lower bounds of what the trained records taught. Probes are graded by numbers appearing in the
answer; a different grader was not compared. The experiments ran on binaries built from the commits
of \cref{sec:repro}, which differ between the first and second wave and, for the repaired-verifier
rerun, from the registered run's; each result names its wave. The training qualification of the 27B
model used random weights for the long sequences and a handful of examples for the real checkpoint;
it shows that steps fit and run, not that the adapter changes generation, and no sequence longer
than 16384 tokens was run.

The cause of the generator's missing subjects is localised, not resolved: the rule is not at fault
beyond the markup case, the generator at two sizes names a subject its question omits, and a changed
brief was not tested. Corpus expansion beyond the first pass (variants, retention queue,
rehearsal, claims) is described from the code; variants and rehearsal were exercised as stage
counts and as a post hoc setting, and none of the mechanisms was evaluated for the value it adds.

\paragraph{The coding-task experiment.} It is not measured. Everything said about it in
\cref{sec:protocol,sec:pending,sec:replicate} is a registered rule or a statement about code; the
sign-test table is a design calculation. The registration fixes one model size for the first
training result (Qwen3-8B, because training, saving and serving an adapter work end to end for
it), one agent harness, and runs with process isolation \code{none} on the node, so the exposure
audit detects but cannot prevent a solver reading hidden files (\cref{sec:audit}).\cl{AU09} Mutation-seeded
tasks test whether a model can undo a synthetic fault, not whether it can debug a historical one;
tasks from repository history are registered as a source, not yet built.\cl{PR10} A family is the file a
fault was seeded in for the authored repositories, so clusters are files, not subsystems.\cl{PR02}
An exam of one repository may hold too few non-tied families to decide, and the registration
labels such a comparison a pilot instead of changing the rule. The base model may have seen
public repositories in pretraining; the paired base arm bounds that effect, it does not remove it.
Agent rollouts of realistic length make training cost unmeasured for the 8B model; the only
training rate measured here is that of the 27B architecture (\cref{sec:throughput}), on one card,
one revision and one sequence length.

\paragraph{Open defects.} The retention check is blind to taught facts, the default anchor bound
refuses on one net lost answer in a small suite, a workspace pass is an exit status, and the
practice gate has no canaries (\cref{sec:open}). The first campaign was not rerun with the
repairs that exist, so its reported outcome describes the build of its time.

\input{generated/limitations}
